\documentclass[tikz,border=18pt]{standalone}
\usepackage{luatexja}
\usepackage{tikz}
\usetikzlibrary{shapes.geometric,arrows.meta,positioning,calc}

\definecolor{ciablue}{RGB}{20, 80, 150}
\definecolor{ciagreen}{RGB}{30, 135, 55}
\definecolor{ciaorange}{RGB}{215, 100, 25}
\definecolor{corebg}{RGB}{245, 248, 255}

\begin{document}
\begin{tikzpicture}[>=Stealth, font=\small]

  % -------------------------------------------------------------
  % 左側：CIAの3要素（情報セキュリティの三位一体）
  % -------------------------------------------------------------
  \node[above, font=\large\bfseries, text=ciablue] at (-5.0, 5.2) {【情報セキュリティの3大要素（CIA）】};

  % 三角形配置の3大要素
  \node[draw=ciablue, fill=blue!10, rounded corners=8pt, line width=1.5pt, minimum width=4.6cm, minimum height=2.2cm, align=center] (c) at (-5.0, 3.8) {
    \textbf{\large Ｃ：機密性}\\[3pt]
    \footnotesize (Confidentiality)\\[3pt]
    \scriptsize 許可された人だけが情報にアクセスできる\\（暗号化、アクセス権限、認証）
  };

  \node[draw=ciagreen, fill=green!10, rounded corners=8pt, line width=1.5pt, minimum width=4.6cm, minimum height=2.2cm, align=center] (i) at (-7.5, 0.4) {
    \textbf{\large Ｉ：完全性}\\[3pt]
    \footnotesize (Integrity)\\[3pt]
    \scriptsize 情報が改ざん・破壊されず正確である\\（ハッシュ値、電子署名、バックアップ）
  };

  \node[draw=ciaorange, fill=orange!10, rounded corners=8pt, line width=1.5pt, minimum width=4.6cm, minimum height=2.2cm, align=center] (a) at (-2.5, 0.4) {
    \textbf{\large Ａ：可用性}\\[3pt]
    \footnotesize (Availability)\\[3pt]
    \scriptsize 必要な時にいつでも情報を使える\\（二重化、DDoS対策、耐障害性）
  };

  \draw[<->, very thick, gray!70] (c.south west) -- (i.north);
  \draw[<->, very thick, gray!70] (c.south east) -- (a.north);
  \draw[<->, very thick, gray!70] (i.east) -- (a.west);

  \node[below, font=\scriptsize, align=center] at (-5.0, -1.0) {
    どれか1つが欠けてもシステムは崩壊する！\\[2pt]
    【例】DDoS攻撃＝可用性破壊、改ざん＝完全性破壊、流出＝機密性破壊
  };

  % -------------------------------------------------------------
  % 右側：多層防御（Defense in Depth）同心円モデル
  % -------------------------------------------------------------
  \node[above, font=\large\bfseries, text=ciablue] at (5.5, 5.2) {【多層防御（Defense in Depth）】};

  % 外側から内側への円
  \fill[blue!8] (5.5, 1.8) circle (3.2cm);
  \draw[blue!60, very thick] (5.5, 1.8) circle (3.2cm);
  \node[above, font=\footnotesize\bfseries, text=blue!80] at (5.5, 4.4) {第1層：入口対策（ファイアウォール・WAF・フィルタ）};

  \fill[blue!18] (5.5, 1.8) circle (2.2cm);
  \draw[blue!80, very thick] (5.5, 1.8) circle (2.2cm);
  \node[above, font=\footnotesize\bfseries, text=blue!90] at (5.5, 3.4) {第2層：内部対策（EDR・最小権限・セグメント）};

  \fill[blue!30] (5.5, 1.8) circle (1.2cm);
  \draw[ciablue, line width=1.5pt] (5.5, 1.8) circle (1.2cm);
  \node[above, font=\footnotesize\bfseries, text=ciablue] at (5.5, 2.4) {第3層：出口対策};

  \node[draw=red!80, fill=white, rounded corners=4pt, thick, align=center, inner sep=3pt] at (5.5, 1.5) {
    \textbf{\textcolor{red!80}{重要データ}}\\[1pt]
    \tiny (暗号化保管)
  };

  \node[below, font=\scriptsize, align=center] at (5.5, -1.0) {
    「1枚の壁」が突破されても次で止める！\\[2pt]
    入口・内部・出口の三位一体で重要資産を護り抜く
  };

\end{tikzpicture}
\end{document}
